\section{Coefficient comparisons and state-valued bars}
\label{sec:ordered-state-bars}

\subsection{State coefficients and multiplication faces}
\label{sb:sec:state-coefficient-bars}

A collision expansion determines state coefficients.  To use one
of them as a bar face, it must satisfy the multiplication identity
required by the bar differential.

Let \(k\) be a field and \(V\) a \(k\)-vector space.
Write \(V((t))\) for formal Laurent series in a variable \(t\)
with coefficients in \(V\) and a finite lower bound on the powers
of \(t\).  A formal state field is a bilinear rule, written as a
linear map \(Y:V\otimes_kV\to V((t))\); write its value as \(Y(a,t)b\).
For \(a,b\in V\), define the normal product
\(a\cdot b=[t^0]Y(a,t)b\), where \([t^0]\) extracts the
constant coefficient.

For the initial test put \(V\) in degree zero.  Denote a tensor
\(a_1\otimes\cdots\otimes a_n\) by the word
\([a_1|\cdots|a_n]\).  On words of lengths one, two and three,
define a length-lowering operator \(\beta\) by
\[
 \beta[a]=0,\qquad
 \beta[a|b]=-[a\cdot b],\qquad
 \beta[a|b|c]=-[a\cdot b|c]+[a|b\cdot c].
\]
Applying it twice gives
\[
 \beta^2[a|b|c]=[(a\cdot b)\cdot c-a\cdot(b\cdot c)].
\]
Thus associativity is necessary before this operator can be a
differential.  The oscillator calculation below gives a nonzero
associator in every characteristic, including characteristic two.

There is a separate comparison between two ways of recording the
coefficients of an order interval.  The first uses polynomial
differential forms with regular endpoint values.  The second records
two endpoint values and one overlap value.  Integration compares these
complexes in characteristic zero but does not preserve multiplication.
We construct its higher corrections on words.  A cut splits a word
into two consecutive subwords; the resulting map preserves these cuts
and the word differential.  With their cut coproducts the bars are
differential graded coalgebras.  No multiplication on the bars is
assumed.  Identifying such a bar with one in a geometric category
requires additional maps, including the appropriate module ends or
augmentation.

We retain the cohomological signs, exterior-square convention and
finite-polynomial state domain of Section~\ref{sec:free-states}.
The field is arbitrary except where polynomial integration is used.
Every formal expansion below has the specified algebraic domain.

\subsection{The two-chart coefficient complex}
\label{sb:sec:two-chart}

Let \(R=k[u]\), \(L=k[u,u^{-1}]\).  The two order charts have
coordinate ring \(R\), and their overlap has ring \(L\).
The coordinate \(s\) records passage from the first chart to the second.
Polynomial forms with regular endpoints form the differential graded algebra
\[
 C^0=\{f(s)\in L[s]:f(0),f(1)\in R\},\qquad
 C^1=L[s]ds,\qquad df=f'(s)ds.
\]
There are no other degrees.  Multiplication is ordinary multiplication
of forms, so \((ds)^2=0\).

For later changes of the spatial base, define the interval algebra
\(\mathcal I_T(v)\) for a commutative ring \(T\) and a
nonzerodivisor \(v\in T\):
\begin{equation}\label{sb:eq:general-interval}
 \mathcal I_T(v)^0=Te_0\oplus Te_1,\qquad
 \mathcal I_T(v)^1=T[v^{-1}]\theta_v,\qquad
 de_0=-\theta_v,\quad de_1=\theta_v.
\end{equation}
There are no other degrees.  The two endpoint actions below define
its multiplication; its unit is \(e_0+e_1\), and its augmentation
is evaluation at the first endpoint.  In the present comparison write
\(A=\mathcal I_R(u)\) and \(\theta=\theta_u\).  Thus
\begin{equation}\label{sb:eq:two-chart}
 A^0=R\oplus R,\qquad A^1=L\theta,\qquad
 d(a_0,a_1)=(a_1-a_0)\theta.
\end{equation}
Its product is
\[
 \begin{aligned}
 (a_0,a_1)(b_0,b_1)&=(a_0b_0,a_1b_1),\\
 (a_0,a_1)(\ell\theta)&=a_0\ell\theta,&
 (\ell\theta)(b_0,b_1)&=\ell b_1\theta,&
 (\ell\theta)(m\theta)&=0.
 \end{aligned}
\]
The endpoint actions prove associativity.  The product rule follows from
\[
 a_1b_1-a_0b_0=(a_1-a_0)b_1+a_0(b_1-b_0).
\]
Products of two degree-one elements vanish, so this is the only
nonzero product-rule calculation.  The unit is \((1,1)\).
The algebra need not be commutative.

Both algebras are augmented over the spatial ring \(R\):
\(\varepsilon_C(f)=f(0)\) and
\(\varepsilon_A(a_0,a_1)=a_0\), with degree-one elements
sent to zero.  These augmentations concern coefficient multiplication.
They do not assert that a vacuum projection preserves collision operations.

\begin{proposition}[The coefficient comparison and its characteristic]
\label{sb:prop:coefficient-comparison}
Over every field,
\(H^0(A)=R\) and \(H^1(A)=L/R\).
In characteristic zero, the map
\[
 I(f)=(f(0),f(1)),\qquad
 I(g(s)ds)=\left(\int_0^1g(s)ds\right)\theta
\]
is an augmentation-preserving quasi-isomorphism of complexes.
In characteristic \(p>0\), the class \(u^{-1}ds\) is exact
in \(C\), while \([u^{-1}\theta]\ne0\) in \(H^1(A)\).
Consequently no chain map in that characteristic can preserve this pole class.
The element \(u^{-1}(s^{p^2}-s^p)\) is also a nonzero closed
degree-zero element of \(C\) with zero endpoint values.
\end{proposition}

\begin{proof}
The kernel of \((a_0,a_1)\mapsto a_1-a_0\) is the diagonal
copy of \(R\); its image is \(R\subset L\).  This proves
the asserted cohomology and the nonzero pole class.

In characteristic zero every positive integer is invertible in \(k\).
Integration means polynomial antiderivation, performed in \(L[s]\).
The identity \(\int_0^1 f'=f(1)-f(0)\) proves the chain equation.
Define a reverse chain map and a degree-minus-one map by
\[
 \begin{aligned}
 W(a_0,a_1)&=(1-s)a_0+sa_1,& W(\ell\theta)&=\ell ds,\\
 h(f)&=0,&
 h(g(s)ds)&=\int_0^s g(t)dt-s\int_0^1g(t)dt.
 \end{aligned}
\]
The last polynomial vanishes at both endpoints, hence belongs to
\(C^0\).  Direct differentiation gives
\(IW=\id\) and \(dh+hd=\id-WI\).
On a function the latter identity is
\(f(s)-f(0)-s(f(1)-f(0))\); on a one-form it is
\((g(s)-\int_0^1g)ds\).  This proves the quasi-isomorphism.

In characteristic \(p\), the polynomial \(s-s^p\) vanishes
at both endpoints and has derivative one.  Thus
\begin{equation}\label{sb:eq:exact-pole}
 u^{-1}ds=d\bigl(u^{-1}(s-s^p)\bigr).
\end{equation}
A chain map takes boundaries to boundaries, proving the obstruction.
Both exponents in \(s^{p^2}-s^p\) are divisible by \(p\), so
its derivative is zero.  Its two monomials are distinct and its
endpoint values vanish, proving the final assertion.
\end{proof}

The map \(I\) is not multiplicative, even in characteristic zero.
For example \(I(s\,ds)=\tfrac12\theta\), but
\(I(s)I(ds)=(0,1)\theta=0\).  The next construction retains
this defect and supplies its higher chain corrections.

\subsection{Suspensions, multiplication signs, and cuts}

For an augmented differential graded algebra \(D\) over a spatial
ring \(R\), write
\(\bar D=\ker\varepsilon_D\).  Let \(s\bar D\) have
\(|sa|=|a|-1\), and put
\[
 B_R(D)=\bigoplus_{n\ge0}(s\bar D)^{\otimes_R n}.
\]
Here a word \([a_1|\cdots|a_n]\) has total degree
\(\sum|a_i|-n\).  Equivalently, the independent bar homological
degree is \(n\), and the total degree is internal degree minus
bar degree.  All words have finite length.  We use the coderivation
determined by
\begin{equation}\label{sb:eq:bar-signs}
 b_1[a]=[da],\qquad
 b_2[a|b]=(-1)^{|a|+1}[ab].
\end{equation}
In a longer word, a degree-one operator applied after the first
\(j-1\) entries carries sign
\((-1)^{\sum_{i<j}(|a_i|-1)}\).
Thus, for degree-zero inputs,
\[
 b[a|b|c]=-[ab|c]+[a|bc].
\]
The augmentation makes the outer faces vanish on these reduced words.
The ground ring is retained under spatial substitutions: a map
\(R\to R'\) must be accompanied by an augmentation-preserving
map to the corresponding \(R'\)-algebra, with its explicit scalar extension.

\begin{lemma}[The bar differential and its cuts]\label{sb:lem:bar}
The operator \(b=b_1+b_2\) has degree one and satisfies \(b^2=0\).
The map
\[
 \Delta_{\mathrm{dec}}[a_1|\cdots|a_n]
 =\sum_{j=0}^n[a_1|\cdots|a_j]\otimes[a_{j+1}|\cdots|a_n]
\]
is coassociative and commutes with \(b\), using the tensor differential.
\end{lemma}

\begin{proof}
The square of \(b_1\) vanishes because \(d^2=0\).
On two entries, the terms in \(b_1b_2+b_2b_1\) are
\[
 (-1)^{|a|+1}[da\,b]-[a\,db]
 +(-1)^{|a|}[da\,b]+[a\,db]=0.
\]
On three entries, the two terms of \(b_2^2\) are
\((-1)^{|b|}[(ab)c]\) and
\(-(-1)^{|b|}[a(bc)]\); they cancel by associativity.
Terms on disjoint positions cancel because the two local operators
have degree one.  These cases exhaust \(b^2\).

Two successive cuts choose three consecutive subwords, independent
of which cut is made first.  This proves coassociativity.
For the chain equation, a differential or multiplication wholly on
one side of a cut occurs on that tensor factor.  Its prefix sign is
the same as in the original word, including the total degree of the
first factor when the operation is on the second.  A multiplication
across a proposed cut leaves no cut between its merged entries and
therefore contributes on neither side.  This accounts for every term.
\end{proof}

If the coefficients of the \(j\)-th input have pole bounds
\(N^{(j)}_\delta\) on a finite list of spatial divisors, multiplication
has bound \(N^{(j)}_\delta+N^{(j+1)}_\delta\).
The linear span of tensors with
\(\sum_jN^{(j)}_\delta\le P_\delta\) is therefore a subcomplex.
The span is part of this definition.  A nonlinear condition on the
sum of separately measured divisorial orders need not define a vector space.

\subsection{A comparison on every bar word}

For this section assume \(\operatorname{char}k=0\).
For polynomials \(g_i\in L[s]\), define
\begin{equation}\label{sb:eq:iterated-integral}
 J_n(g_1,\ldots,g_n)=
 \int_{0\le s_1\le\cdots\le s_n\le1}
       g_1(s_1)\cdots g_n(s_n)\,ds_1\cdots ds_n.
\end{equation}
This is an algebraic instruction: integrate \(s_1\) from zero to
\(s_2\), then \(s_2\) to \(s_3\), and so on.
For a monomial it is explicitly
\[
 J_n(s^{r_1},\ldots,s^{r_n})
 =\prod_{j=1}^n(r_1+\cdots+r_j+j)^{-1}.
\]
The formula follows by successive antiderivatives and defines the
integral over any characteristic-zero field.

Define maps from words to length-one words by
\[
 \mathcal I_1[a]=[I(a)],\qquad
 \mathcal I_n[g_1ds|\cdots|g_nds]
       =[J_n(g_1,\ldots,g_n)\theta]\quad(n>1),
\]
and set \(\mathcal I_n=0\) for \(n>1\) if one input is a
degree-zero function.  Extend linearly.  A word of one-forms has total
degree zero, as does its length-one image, so these maps have degree zero.

The map on full bars sums over consecutive partitions:
\begin{equation}\label{sb:eq:bar-map}
 \mathcal I[a_1|\cdots|a_n]
 =\sum_{0=i_0<i_1<\cdots<i_r=n}
   [\mathcal I_{i_1-i_0}(a_1,\ldots,a_{i_1})|\cdots|
    \mathcal I_{i_r-i_{r-1}}(a_{i_{r-1}+1},\ldots,a_n)].
\end{equation}
Inside the brackets on the right, \(\mathcal I_m(\cdots)\)
means its underlying element of \(\bar A\), without an extra suspension.
The empty word maps to the empty word.  The sum is finite.

Write \(b_C\) and \(b_A\) for the differentials of
\(B_R(C)\) and \(B_R(A)\), defined by \eqref{sb:eq:bar-signs}.
On either bar let \(\Delta_{\mathrm{dec}}\) be the cut coproduct
of Lemma~\ref{sb:lem:bar}.  The two required identities are
\begin{equation}\label{sb:eq:spine}
 b_A\mathcal I=\mathcal I b_C,
 \qquad
 \Delta_{\mathrm{dec}}\mathcal I
       =(\mathcal I\otimes\mathcal I)\Delta_{\mathrm{dec}}.
\end{equation}

\begin{theorem}[The two-chart bar comparison]\label{sb:thm:bar-comparison}
Over a field of characteristic zero, formula~\eqref{sb:eq:bar-map}
is a quasi-isomorphism of coaugmented
differential graded \(R\)-coalgebras \(B_R(C)\to B_R(A)\).
It preserves the filtration by maximal word length and the total
pole tuple.  Its component from length \(n\) to length one has
pole bound \(\sum_iN^{(i)}_\delta\) on each divisor.
\end{theorem}

\begin{proof}
Cutting a consecutive partition at a boundary between its blocks is
the same as first cutting the original word and then partitioning
the two resulting words.  This bijection proves the coalgebra equation
in \eqref{sb:eq:spine}, including the empty subwords.

It remains to check the chain equation after projection to length one.
That check suffices: the difference \(b_A\mathcal I-\mathcal I b_C\)
is a coderivation relative to \(\mathcal I\).  If its length-one
projection vanishes, induction on the input length makes every proper
cut of its value zero.  A tensor-coalgebra element with every proper
cut zero has length one.  Its zero projection then makes it zero.

For a word of one-forms, all differentials and all adjacent products
vanish in both coefficient algebras.  For a word containing at least
two functions and of length greater than two, every relevant higher
component is zero: differentiating or multiplying once leaves a function
in a block of length greater than one.  On two functions the remaining
identity is endpoint evaluation of their product.  Length one is
Proposition~\ref{sb:prop:coefficient-comparison}.

Consider a word with exactly one function \(f\), in position \(j\),
and one-forms elsewhere.  Every preceding suspended one-form has degree
zero.  In \(\mathcal I b_C\), differentiating \(f\) gives a
positive \(J_n(\ldots,f',\ldots)\).  Merging the preceding
one-form with \(f\) gives a positive term; merging \(f\) with the
following one-form gives a negative term, by \eqref{sb:eq:bar-signs}.
For \(1<j<n\), integration in its one variable gives
\[
 \int_{s_{j-1}}^{s_{j+1}}f'(s_j)ds_j
       =f(s_{j+1})-f(s_{j-1}).
\]
This cancels the two multiplication terms.  The length-one projection
of \(b_A\mathcal I\) is also zero, because a nonzero higher block
cannot contain \(f\), and an interior singleton cannot be an endpoint
of a partition into two blocks.

If \(j=1\), the same computation leaves
\(-f(0)J_{n-1}(g_2,\ldots,g_n)\theta\), exactly the left
endpoint action with the negative multiplication sign in \(A\).
If \(j=n\), it leaves
\(J_{n-1}(g_1,\ldots,g_{n-1})f(1)\theta\), exactly the right
endpoint action with its positive sign.  These cases prove the chain
equation in all lengths.

Filter both bars by words of length at most \(n\).
On the length-\(n\) associated quotient, the map is
\((sI)^{\otimes n}\), with only the internal differential.
The explicit \(R\)-linear contraction in Proposition~\ref{sb:prop:coefficient-comparison}
restricts to the augmentation ideals and proves that \(I\) is a
chain homotopy equivalence there.  Tensoring a chain homotopy
equivalence with any complex preserves its homotopy inverse: tensor
the two inverse maps and the homotopy with the identity, using the
graded tensor signs.  Applying this one factor at a time proves
that \((sI)^{\otimes_R n}\) is a chain homotopy equivalence.
Induction on the finite word filtration gives a quasi-isomorphism
on every filtered piece.  Taking their union preserves kernels modulo
images, since each element and each equation involve finite words.
This proves the assertion for the direct-sum bars.

Integration changes only order coordinates.  Every integrand in
\(J_n\) is a product of the original spatial coefficients; it has
pole bound \(\sum_iN^{(i)}_\delta\).  Antiderivation multiplies
by scalars of \(k\) and adds such coefficients.  It introduces no
new spatial divisor.  The same argument applies separately to every
block in \eqref{sb:eq:bar-map}, proving both filtration statements.
\end{proof}

\begin{example}[Three poles and all cuts]\label{sb:ex:three-pole-bar}
Put \(c=ds/u\) and \(e_N=\theta/u^N\).  Then
\[
 \begin{aligned}
 \mathcal I[c|c]&=[e_1|e_1]+\tfrac12[e_2],\\
 \mathcal I[c|c|c]&=[e_1|e_1|e_1]
    +\tfrac12[e_2|e_1]+\tfrac12[e_1|e_2]+\tfrac16[e_3].
 \end{aligned}
\]
All displayed elements are cycles.  The proper cuts of the second
line are
\[
 [e_1]\otimes[e_1|e_1]
 +[e_1|e_1]\otimes[e_1]
 +\tfrac12[e_2]\otimes[e_1]
 +\tfrac12[e_1]\otimes[e_2],
\]
which equal \((\mathcal I\otimes\mathcal I)\) applied to the
two proper cuts of \([c|c|c]\).  Each summand has total pole
bound three.  The length-one correction has exact pole order three.
The fractions \(1/2\) and \(1/6\) belong to the coefficient
comparison; the free-state Hasse normalization above uses no such divisions.
\end{example}

\subsection{State coefficients and the binary operation}

First assume that \(k\) has characteristic zero, and let \(D\)
be any augmented differential graded \(R\)-algebra.
The construction above extends to \(C\otimes_R D\) and
\(A\otimes_R D\), with coefficient factors first.  Its linear
part is \(I\otimes\id_D\).  For \(n>1\), its higher
component is zero unless every \(C\)-factor is a one-form, and then
\begin{equation}\label{sb:eq:extension}
 [g_1ds\otimes v_1|\cdots|g_nds\otimes v_n]
 \longmapsto
 [J_n(g_1,\ldots,g_n)\theta\otimes v_1\cdots v_n].
\end{equation}
There is no extra sign in this formula.  A suspended one-form with
state \(v_i\) has degree \(|v_i|\).

\begin{lemma}[Extension to state coefficients]\label{sb:lem:state-extension}
Over a field of characteristic zero, formula~\eqref{sb:eq:extension},
summed over consecutive partitions,
is an augmentation-preserving coalgebra chain quasi-isomorphism
\(B_R(C\otimes_R D)\to B_R(A\otimes_R D)\).
It is natural in strict maps of augmented differential graded algebras
\(D\) and in augmentation-preserving spatial coefficient substitutions,
after the stated scalar extension, for which all displayed expansions
are defined.
\end{lemma}

\begin{proof}
For a word of one-form coefficients, differentiating its state product
on the right gives sign \(-(-1)^{\sum_{i<j}|v_i|}\) on
\(dv_j\).  The minus is from the leading \(\theta\).
The same sign occurs on the left: the prefix suspended degree is
\(\sum_{i<j}|v_i|\), and
\(d(g_jds\otimes v_j)=-g_jds\otimes dv_j\).
These terms agree.

For a word with one function coefficient in position \(j\),
the three terms in the one-variable integration argument have the
common sign \((-1)^{\sum_{i<j}|v_i|}\).
Indeed, the multiplication of that function-state with the next
one-form-state contributes both the bar sign and the sign for moving
\(v_j\) across \(ds\); together these give the required minus.
The preceding multiplication gives the required plus.  State products
stay in their original order, so associativity of \(D\) suffices;
commutativity of \(D\) is not needed.  Terms differentiating
\(D\)-factors when a function coefficient remains have zero higher
component.  With two or more function coefficients the same vanishing
argument and the length-two algebra identity complete the check.

The consecutive-partition proof gives the coalgebra identity.
The filtration proof of Theorem~\ref{sb:thm:bar-comparison} applies
to \(I\otimes\id_D\), since tensoring the explicit coefficient
contraction over \(R\) preserves its chain homotopy.
All formulas commute with the stated maps by linearity, multiplication,
and coefficientwise antiderivation.
\end{proof}

\begin{proposition}[The length-one obstruction in positive characteristic]
\label{sb:prop:length-one-obstruction}
Suppose \(\operatorname{char}k=p>0\).  The endpoint rule
\(I(f)=(f(0),f(1))\) on functions and the constant-form rule
\(I(\ell ds)=\ell\theta\), \(\ell\in L\), cannot be the
length-one component of a coaugmented coalgebra chain map
\(B_R(C)\to B_R(A)\).  Consequently the extension of
Lemma~\ref{sb:lem:state-extension} does not hold over arbitrary fields,
even for \(D=R\).
\end{proposition}

\begin{proof}
The polynomial
\[
 f(s)=\frac{s-s^p}{u}
\]
has zero endpoint values, so it belongs to the augmentation ideal
of \(C^0\).  Its differential is \(df=ds/u\).
The two prescribed values would therefore be
\[
 I(f)=0,\qquad I(df)=\frac{\theta}{u}.
\]
If \(\mathcal I\) were such a map, these prescribed values on
the length-one word \([f]\) would give
\[
 b_A[I(f)]=0,
 \qquad
 \mathcal I b_C[f]=[I(df)]=[\theta/u]\ne0.
\]
The last inequality is a statement about chains: \(s\bar A\)
is a direct summand of \(B_R(A)\), and \(\theta/u\ne0\)
in \(A^1\).  A coaugmented coalgebra map sends a length-one word,
which has no proper cut, to a length-one word.  Higher components
act on at least two inputs and cannot change this failure.
Taking \(D=R\) leaves the same coefficient complex and proves
the assertion about state extension.  The rational normalization
\(J_2(1,1)=1/2\) also cannot be specialized to characteristic two.
\end{proof}

For the direct binary operation, take the one-species instance of
the free state algebra \(S\) in \eqref{fs:eq:state-algebra} over any
field.  Use its auxiliary weight, Hasse translation and kernels from
\eqref{fs:eq:hasse-translation}--\eqref{fs:eq:kernel-translation}.
Write \(P_{rs}=P_{12}^{rs}\) for the operator in
\eqref{fs:eq:pair-operators}.  All sums act on finite polynomial inputs.
Write \(E_u(a\otimes b)=(\tau_u a)b\).  The binary form operation is
\[
 Y_C=E_u\left(1+\sum_{r,s}ds\,K_{rs}(u)\otimes P_{rs}\right).
\]
The sum is finite on a polynomial input.  Its values lie in the
product over state weights, with one finite common spatial pole bound.
Its coefficient differential is zero, and each contraction has total
degree zero.

\begin{proposition}[The binary operation on the actual overlap complex]
\label{sb:prop:binary-transfer}
Over every field the formula
\[
 Y_A(a,b)=((\tau_u a)b,(\tau_u a)b)
    +\sum_{r,s}\theta K_{rs}(u)\otimes E_uP_{rs}(a\otimes b)
\]
is a degree-zero chain map to the weight-completed \(A\otimes S\).
In characteristic zero it equals \((I\otimes\id)Y_C\).
Its pole bounds are exactly those of its displayed contraction terms,
subject to cancellation or vanishing of their integer coefficients.
\end{proposition}

\begin{proof}
The uncontracted coefficient has equal endpoint values and is closed
in \(A\).  Each contracted coefficient lies in \(A^1\), whose
differential is zero.  The state differential is zero.  The form or
overlap generator has degree one and the derivative has degree minus
one, proving the total degree assertion.  The equality in characteristic
zero is integration of a constant one-form.  Translation has polynomial
coefficients in each fixed state weight and adds no spatial poles.
\end{proof}

For instance
\[
 Y_A(x_0,\eta_0)=((\tau_u x_0)\eta_0,(\tau_u x_0)\eta_0)
                         +\frac{\theta}{u}\,1.
\]
The class of its scalar-state term is nonzero in every characteristic.
The state augmentation \(S\to k\) kills both inputs but sends
that output to \(\theta/u\).  Therefore this augmentation cannot
preserve the binary collision operation, even up to a chain homotopy.
This obstruction uses \(A\); using \(C\) in characteristic
\(p\) would erase it by \eqref{sb:eq:exact-pole}.

When \(k\) has characteristic zero,
Lemma~\ref{sb:lem:state-extension} applies with state algebra
\(D=R\otimes_k S\).  This bar comparison concerns the
ordinary multiplication of \(C\otimes S\) and \(A\otimes S\).
Proposition~\ref{sb:prop:binary-transfer} separately identifies their binary
collision maps.  It does not make a collision map into that multiplication.

\subsection{Heisenberg fields with integral jets}

Fix \(\kappa\in k\) and let
\(V=k[h_0,h_1,h_2,\ldots]\), in cohomological degree zero,
with vacuum \(1\) and differential zero.  Define creation and
annihilation operators by
\[
 a_{-m}=h_{m-1}\cdot,\qquad
 a_m=\kappa m\partial_{h_{m-1}}\quad(m\ge1),\qquad a_0=0.
\]
They satisfy
\([a_m,a_n]=\kappa m\delta_{m,-n}\id\).
This follows by differentiating one polynomial variable; all other
operator pairs commute.  The field and its two parts are
\[
 h(z)=H(z)+h_+(z),\qquad
 H(z)=\sum_{r\ge0}h_rz^r,\qquad
 h_+(z)=\kappa\sum_{m\ge1}m\partial_{h_{m-1}}z^{-m-1}.
\]
Every polynomial state is killed by all sufficiently large annihilation
modes.  The auxiliary weight is \(\operatorname{wt}(h_r)=r+1\); fixed
weight spaces are finite-dimensional.

For an integer \(r\ge0\), let
\(\partial_z^{[r]}z^n=\binom nr z^{n-r}\), also for negative
integers \(n\).  Put
\(h^{[r]}(z)=\partial_z^{[r]}h(z)\).
For a polynomial state \(h_{r_1}\cdots h_{r_m}\), its field
is the normally ordered product of these jet fields: expand each
field into creation and annihilation parts and move every creation
operator to the left without adding commutators.  Operators within
each part commute, so this prescription is independent of the order
of the factors.  It uses integer coefficients only.
On the vacuum this field gives \(\tau_z(h_{r_1}\cdots h_{r_m})\), where
\[
 \tau_z h_r=\sum_{m\ge0}\binom{r+m}{r}h_{r+m}z^m.
\]
We shall use only the displayed fields of one and two generators and their
action on a one-generator state; these operators are already completely defined.

The basic commutator is
\begin{equation}\label{sb:eq:heisenberg-contraction}
 [h_+(z),H(w)]
 =\kappa\sum_{m\ge1}m w^{m-1}z^{-m-1}
 =\kappa\,\iota_{z,w}(z-w)^{-2}.
\end{equation}
Here \(\iota_{z,w}\) expands in nonnegative powers of \(w/z\).
The second equality follows by multiplying the geometric series by
itself, whose coefficient of degree \(m-1\) is \(m\).
It therefore holds over every field.

\begin{lemma}[Jet contractions]\label{sb:lem:heisenberg-jets}
The contraction of \(h_r\) with \(h_s\) has rational coefficient
\begin{equation}\label{sb:eq:heisenberg-kernel}
 B_{rs}(u)=\kappa(-1)^r(r+s+1)\binom{r+s}{r}u^{-r-s-2}.
\end{equation}
Moreover,
\[
 Y(h_rh_s,z)h_t
 =\tau_z(h_rh_s)h_t
   +B_{rt}(z)\tau_z h_s+B_{st}(z)\tau_z h_r.
\]
\end{lemma}

\begin{proof}
The coefficient of \(v^rw^s\) in \((u+v-w)^{-2}\) is
\((-1)^r(r+s+1)\binom{r+s}{r}u^{-r-s-2}\): choose the
\(r\) positive and \(s\) negative displacements in the
geometric-square expansion.  Applying Hasse derivatives to
\eqref{sb:eq:heisenberg-contraction} gives this coefficient.
In the normally ordered two-generator field acting on \(h_t\),
the term with two annihilators vanishes.  Either one annihilator
contracts with \(h_t\), leaving the other creation field, or both
factors are creation fields.  These are exactly the three displayed terms.
\end{proof}

\subsection{Three-input collisions and their exact poles}

Apply \(h(z_1)h(z_2)h(z_3)\) to the vacuum in the ordered
expansion domain.  Moving the annihilators to the right using
\eqref{sb:eq:heisenberg-contraction} gives the rational representative
\[
 \mathcal H_3(z_1,z_2,z_3)
 =H(z_1)H(z_2)H(z_3)
  +\sum_{i<j}\frac{\kappa H(z_k)}{(z_i-z_j)^2},
 \quad \{k\}=\{1,2,3\}\setminus\{i,j\}.
\]
There can be at most one contraction among three one-generator inputs.
Thus the commutator calculation proves this identity, with no infinite
sum of contractions or factorial normalization.

Set \(z_1=x+t,z_2=x,z_3=0\), and work first in
\[
 Q=k[t,x,t^{-1},x^{-1},(x+t)^{-1}].
\]
The two collision expansions have different targets:
\[
 \begin{aligned}
 \iota_L:Q&\longrightarrow k[x,x^{-1}]((t)),&
 (x+t)^{-2}&\longmapsto\sum_{m\ge0}(-1)^m(m+1)x^{-m-2}t^m,\\
 \iota_R:Q&\longrightarrow k[u,u^{-1}]((s)),&
 x&\longmapsto s,\quad t\longmapsto u-s.
 \end{aligned}
\]
For \(\iota_R\), expand \((u-s)^{-2}\) in nonnegative
powers of \(s/u\).  These maps are injective: each sends the
specified denominators to units, and a polynomial mapping to zero
would have every coefficient zero.  They do not identify their two
Laurent targets.  State coefficients are taken weight by weight.

\begin{proposition}[The two three-input expansions]\label{sb:prop:heisenberg-three}
Writing \(h=h_0\), the common rational representative is
\begin{equation}\label{sb:eq:heisenberg-three}
 H(x+t)H(x)h+\frac{\kappa h}{t^2}
       +\frac{\kappa H(x)}{(x+t)^2}
       +\frac{\kappa H(x+t)}{x^2}.
\end{equation}
Its left expansion equals
\(Y(Y(h,t)h,x)h\), with
\[
 Y(h,t)h=\kappa t^{-2}1+\sum_{m\ge0}h_mh\,t^m.
\]
Its right expansion equals \(Y(h,u)Y(h,s)h\), with
\(u=x+t,s=x\).
For every \(m\ge0\), its left coefficient of \(t^m\) is
\begin{equation}\label{sb:eq:heisenberg-coefficient}
 \tau_x(h_mh)h
   +\kappa(-1)^m(m+1)x^{-m-2}H(x)
   +\kappa x^{-2}\tau_xh_m.
\end{equation}
\end{proposition}

\begin{proof}
The rational representative was obtained by the three possible pair
contractions.  On the left, insert the binary expansion just displayed
and apply Lemma~\ref{sb:lem:heisenberg-jets} to \(Y(h_mh,x)h\).
Since \(B_{m0}=\kappa(-1)^m(m+1)x^{-m-2}\) and
\(B_{00}=\kappa x^{-2}\), its coefficients are
\eqref{sb:eq:heisenberg-coefficient}.  The uncontracted sum is
\(H(x+t)H(x)h\), by the Hasse translation law.
The negative inner coefficient contributes \(\kappa t^{-2}h\).
This proves the left equality coefficient by coefficient.

On the right, \(Y(h,s)h=H(s)h+\kappa s^{-2}\).
The creation part of \(Y(h,u)\) gives
\(H(u)H(s)h+\kappa s^{-2}H(u)\).
Its annihilation part gives
\(\kappa\iota_{u,s}(u-s)^{-2}h+\kappa u^{-2}H(s)\).
Substitution of \(u=x+t,s=x\) yields exactly the right expansion
of \eqref{sb:eq:heisenberg-three}.
\end{proof}

Assume \(\kappa\ne0\).  The exact outer pole of
\eqref{sb:eq:heisenberg-coefficient} can be read from its one-generator
state component.  For \(m>0\), if \(m+1\ne0\) in \(k\),
its \(h_0\)-coefficient is
\(\kappa(-1)^m(m+1)x^{-m-2}\), so the exact order is \(m+2\).
If \(m+1=0\) in \(k\), that entire term vanishes, leaving a
pole of exact order two detected by \(h_m\).
For \(m=0\), the singular part is
\(2\kappa x^{-2}H(x)\): its order is two if
\(\operatorname{char}k\ne2\), and the whole coefficient is
regular in characteristic two.  The first cases are
\[
 \begin{array}{c|ccc}
 &m=0&m=1&m=2\\\hline
 \operatorname{char}k=0&2&3&4\\
 \operatorname{char}k=2&0&2&4\\
 \operatorname{char}k=3&2&3&2.
 \end{array}
\]
Order zero means absence of a pole.  These are actual state coefficients;
neither a scalar common-denominator estimate nor a sum of separate
divisorial maxima determines the table.

\subsection{The length-three bar obstruction}

Define \(a\cdot b=[z^0]Y(a,z)b\) on the specified polynomial
states.  For single-generator states, \(h_r\cdot h_s=h_rh_s\).
The product is unital, but the following calculation rules out using
it as an associative bar face.

\begin{theorem}[A characteristic-independent associator]
\label{sb:thm:heisenberg-associator}
Over every field,
\begin{equation}\label{sb:eq:heisenberg-associator}
 (h_0\cdot h_1)\cdot h_0-h_0\cdot(h_1\cdot h_0)
       =\kappa h_3.
\end{equation}
For \(\kappa\ne0\), the proposed reduced bar face
\(b[a|b]=-[a\cdot b]\) has
\(b^2[h_0|h_1|h_0]=[\kappa h_3]\ne0\).
No higher multiplication can make this product into an
\(A_\infty\)-algebra on the same zero-differential state complex.
\end{theorem}

\begin{proof}
Lemma~\ref{sb:lem:heisenberg-jets} gives
\[
 Y(h_0h_1,z)h_0
 =\tau_z(h_0h_1)h_0
  +\kappa z^{-2}\tau_z h_1-2\kappa z^{-3}H(z).
\]
The constant coefficients of the last two terms are
\(3\kappa h_3\) and \(-2\kappa h_3\), respectively.
Their sum is \(\kappa h_3\) over the integers and hence over
every field.  The uncontracted constant coefficient is \(h_0^2h_1\).
In \(Y(h_0,z)(h_1h_0)\), all annihilation terms have strictly
negative powers of \(z\), and the creation constant is
\(h_0^2h_1\).  This proves \eqref{sb:eq:heisenberg-associator}.

Applying the two adjacent faces with \eqref{sb:eq:bar-signs} gives
\[
 b^2[a|b|c]=[(a\cdot b)\cdot c-a\cdot(b\cdot c)].
\]
The resulting nonzero vector proves the failure.  In a bar coderivation
with a higher component from \(m_3\), the length-three,
length-one component of its square is
\(b_2^2+b_1b_3+b_3b_1\).  Here \(b_1=0\), so the associator
would still have to vanish.  No choice of \(m_3\), or of components
of greater arity, can alter this equation.
\end{proof}

For comparison, the unjetted triple has associator
\((h_0\cdot h_0)\cdot h_0-h_0\cdot(h_0\cdot h_0)
=2\kappa h_2\), which vanishes in characteristic two.
The jet in \eqref{sb:eq:heisenberg-associator} is needed to detect the
failure uniformly in the characteristic.

Deconcatenation still makes any binary operation into a coderivation
on tensor words.  Thus checking only that cuts commute with the proposed
face operator would not detect this failure.  A bar complex needs both
the cut identity and \(b^2=0\).

\subsection{The augmentation and the native carrier}

A collision-preserving map of state objects is required to commute
with every coefficient operation
\(a_{(n)}b=[z^{-n-1}]Y(a,z)b\) and to preserve the vacuum.
The one-dimensional vacuum object \(k1\) has
\(Y(1,z)1=1\); all its products with \(n\ge0\) vanish.

\begin{proposition}[No vacuum augmentation at nonzero level]
\label{sb:prop:augmentation-obstruction}
If \(\kappa\ne0\), there is no vacuum-preserving,
collision-preserving linear map \(V\to k1\).
There is also no unital multiplication \(V\otimes V\to V\)
preserving the tensor-product field operations and restricting to
the identity on \(V\otimes1\) and \(1\otimes V\).
These obstructions persist for compatibility up to chain homotopy
on the zero-differential state complexes.
\end{proposition}

\begin{proof}
The binary field gives \((h_0)_{(1)}h_0=\kappa1\).
A proposed augmentation would send its left side to a nonnegative
product of two elements of \(k1\), hence to zero, and its right
side to \(\kappa1\ne0\).

For the second assertion, the tensor-product field satisfies
\[
 Y_{V\otimes V}(h_0\otimes1,z)(1\otimes h_0)
       =H(z)\otimes h_0.
\]
It is regular, so its coefficient of \(z^{-2}\) is zero.
The two unit restrictions send its two inputs to \(h_0,h_0\),
whose coefficient of \(z^{-2}\) is \(\kappa1\).
This contradicts collision preservation.  Chain homotopies induce
equal maps on cohomology.  Since the state differentials are zero,
the same nonzero classes give the same contradictions there.
\end{proof}

The internal bar with two unit ends in a monoidal category
\((\mathcal F,\otimes,\mathbf1)\) starts with an augmented associative algebra object
\(\mathbf1\to A\to\mathbf1\).
Its degree-\(n\) term is \(A^{\otimes n}\); multiplication,
augmentation and unit give its faces and degeneracies.
If \(\mathcal F\) consists of factorization objects, each of
these maps must also preserve their collision structure.
Consequently, any state realization preserving the tensor product,
vacuum and collision coefficients must pass the two tests in
Proposition~\ref{sb:prop:augmentation-obstruction}.
The nonzero-level Heisenberg states with the stated field operations
cannot furnish those data.  This is an obstruction to that specified
realization, not a nonexistence theorem for every geometric bar.

An augmentation of the ordinary polynomial algebra \(V\), obtained
by killing all \(h_r\), is well defined.  Its ordinary associative
bar also exists.  Proposition~\ref{sb:prop:augmentation-obstruction}
proves precisely why that augmentation cannot be substituted for a
collision-preserving one.  Likewise, an ordinary coefficient-bar
comparison does not supply the missing geometric multiplication.

\subsection{A two-sided bar retaining the Heisenberg modes}
\label{sec:mode-bar}

There is an actual associative carrier for the mode operators, with
module ends in place of a character.  Let
\(V=k[h_0,h_1,\ldots]\) be the polynomial Heisenberg Fock space
defined above, with \(\operatorname{wt}(h_j)=j+1\).
Let \(\mathscr W_\kappa\)
be the unital associative \(k\)-algebra on generators
\(p_m,q_m\), \(m\ge1\), with
\[
 [p_m,q_n]=\kappa m\delta_{mn},\qquad
 [p_m,p_n]=[q_m,q_n]=0.
\]
All elements are finite linear combinations of finite words.
The actions \(q_m=h_{m-1}\cdot\) and
\(p_m=\kappa m\partial_{h_{m-1}}\) make \(V\) a left
\(\mathscr W_\kappa\)-module over every field.

Ordered monomials \(q^\alpha p^\beta\), with finite-support
multiindices, form a basis.  The commutation relations first put every
word into that order.  To prove independence, represent the algebra
on \(k[Q_m,P_m:m\ge1]\) by
\(q_m=Q_m\cdot\) and \(p_m=\kappa m\partial_{Q_m}+P_m\cdot\).
These operators satisfy the relations, and
\(q^\alpha p^\beta(1)=Q^\alpha P^\beta\).
Distinct ordered monomials therefore cannot have a linear relation.
This proof also shows that modes remain elements of the algebra when
their action on the particular Fock module vanishes in positive characteristic.

Write \(V=\bigoplus_{E\ge0}V_E\) for its positive weight grading
and define the restricted dual
\(M=\bigoplus_{E\ge0}\operatorname{Hom}_k(V_E,k)\).
The right action is
\((\lambda a)(v)=\lambda(av)\).
A mode shifts weight by a fixed integer, and a finite word has a fixed
total shift.  Thus the right action preserves the restricted dual.
All tensor products in the following complex are algebraic:
\begin{equation}\label{mb:eq:mode-bar}
 \mathcal B_r=M\otimes_k\mathscr W_\kappa^{\otimes_k r}\otimes_k V,
 \qquad |\mathcal B_r|=-r.
\end{equation}

\begin{theorem}[The mode bar and its cuts]\label{mb:thm:mode-bar}
Over every field, the differential
\[
 \begin{aligned}
 d\bigl(\lambda[a_1|\cdots|a_r]v\bigr)
 &=\lambda a_1[a_2|\cdots|a_r]v\\
 &\quad+\sum_{j=1}^{r-1}(-1)^j
   \lambda[a_1|\cdots|a_ja_{j+1}|\cdots|a_r]v\\
 &\quad+(-1)^r\lambda[a_1|\cdots|a_{r-1}]a_rv
 \end{aligned}
\]
squares to zero.  Pairing in length zero and zero in positive lengths
defines a chain map \(\mathcal B_\bullet\to k\).
Its homology is \(\operatorname{Tor}^{\mathscr W_\kappa}_r(M,V)\).
There is a chain cut map
\begin{equation}\label{mb:eq:module-cuts}
 \lambda[a_1|\cdots|a_r]v\longmapsto
 \sum_{i=0}^r
 \lambda[a_1|\cdots|a_i]1
   \otimes_{\mathscr W_\kappa}
 1[a_{i+1}|\cdots|a_r]v
\end{equation}
to the total complex
\(B(M,\mathscr W_\kappa,\mathscr W_\kappa)
 \otimes_{\mathscr W_\kappa}
 B(\mathscr W_\kappa,\mathscr W_\kappa,V)\).
Successive cuts agree after the canonical associativity identification
of the balanced tensor products.
\end{theorem}

\begin{proof}
Denote the unsigned faces by \(d_i\).
For \(i<j\), the identity \(d_id_j=d_{j-1}d_i\)
is associativity when two multiplications meet, the right or left
module identity at an end, and commutation of operations on disjoint
positions otherwise.  In the square of \(\sum(-1)^i d_i\),
these two occurrences have opposite signs.  This proves \(d^2=0\).
Pairing kills the length-one differential because
\((\lambda a)(v)=\lambda(av)\); higher lengths cannot map
nontrivially to length zero in one step.

To identify the homology, use
\(P_r=\mathscr W_\kappa\otimes_k
\mathscr W_\kappa^{\otimes_k r}\otimes_kV\), with the same
faces and augmentation \(P_0\to V\) given by the action.
Each \(P_r\) is a free left \(\mathscr W_\kappa\)-module:
choose a \(k\)-basis for its remaining tensor factors.
The maps \(v\mapsto1[]v\) and
\(a_0[a_1|\cdots|a_r]v\mapsto1[a_0|a_1|\cdots|a_r]v\)
contract the augmented complex as a complex of vector spaces.
Its first face is the original element, and every other face cancels
the corresponding term obtained by first applying the differential.
Thus \(P_\bullet\to V\) is a free module resolution.
Tensoring it on the left with \(M\) gives exactly
\eqref{mb:eq:mode-bar}, proving the Tor assertion.  The contracting map
need not be module-linear; exactness is the property required here.

For the cut map, faces wholly on one side of a cut match the
corresponding tensor differential, with sign \((-1)^i\) when
the differential acts on a second factor following length \(i\).
At the middle, the last face on the left of cut \(i\) gives
\[
 (-1)^i\lambda[a_1|\cdots|a_{i-1}]a_i
            \otimes 1[a_{i+1}|\cdots|a_r]v.
\]
Balancing over \(\mathscr W_\kappa\) moves \(a_i\) to
the initial coefficient of the right factor.  This cancels the first
face on the right of cut \(i-1\), whose sign is \((-1)^{i-1}\).
All middle terms cancel this way.  Two cuts merely choose three
consecutive subwords and insert two copies of the regular module,
so their two orders give the same balanced tensor.
\end{proof}

These cuts pass through the regular module; they do not make
\(M\otimes\mathscr W_\kappa^{\otimes r}\otimes V\)
into a coalgebra by silently identifying its two ends with \(k\).
For \(\kappa\ne0\), the algebra \(\mathscr W_\kappa\)
has no character to \(k\): a character would send
\([p_1,q_1]=\kappa1\) to zero.  The two-sided bar needs no character.

For a concrete calculation, let \(\lambda_n\) take the coefficient
of \(h_0^n\) and kill all other monomials.  Then
\(\lambda_0p_1=\kappa\lambda_1\),
\(\lambda_1q_1=\lambda_0\), and \(p_1h_0=\kappa1\).
The length-two differential is
\[
 d\bigl(\lambda_0[p_1|q_1]1\bigr)
 =\kappa\lambda_1[q_1]1-\lambda_0[p_1q_1]1
       +\lambda_0[p_1]h_0.
\]
Differentiating the middle term gives zero, since
\(\lambda_0p_1q_1=\kappa\lambda_0\) and
\(p_1q_1(1)=\kappa1\).  The other two terms have opposite
differentials
\(\kappa\lambda_0[]1-\kappa\lambda_1[]h_0\).
The central vacuum contribution is retained in the cancellation.
It is not set to zero by an algebra augmentation.

\subsubsection{Evaluation on the Fock-module ends}
\label{mbh:sec}

The derived tensor product in Theorem~\ref{mb:thm:mode-bar} can be
evaluated without changing either module end.  Its differential reduces
to the right action of the annihilation operators on the restricted
dual.  An explicit comparison works over an arbitrary commutative base;
the homology of the resulting complex then depends on the level and
the characteristic.

For the comparison, let \(S\) be a commutative ring and fix
\(c_m\in S\), \(m\ge1\).  Let \(W\) be the unital \(S\)-algebra
with relations
\[
 [p_m,q_n]=c_m\delta_{mn},\qquad
 [p_m,p_n]=[q_m,q_n]=0,
\]
and put \(V_S=S[q_1,q_2,\ldots]\), with
\(q_m=q_m\cdot\) and \(p_m=c_m\partial_{q_m}\).
The preceding mode algebra is obtained by taking \(S=k\) and
\(c_m=\kappa m\), with \(q_m=h_{m-1}\).
All words, polynomials, exterior tensors and sums here have finite
algebraic support.

The ordered-monomial argument above is valid over \(S\).
Indeed, \(q_m=Q_m\cdot\) and
\(p_m=c_m\partial_{Q_m}+P_m\cdot\) act on
\(S[Q_m,P_m:m\ge1]\).  The represented ordered monomial
\(q^\alpha p^\beta\) sends \(1\) to \(Q^\alpha P^\beta\).
Together with reordering, this proves that
the \(q^\alpha p^\beta\) form an \(S\)-basis of \(W\), including
when some \(c_m\) vanish.

Let \(E_S=\bigoplus_{m\ge1}S\epsilon_m\), with each
\(\epsilon_m\) in homological degree one.  For a finite ordered set
\(I=\{i_1<\cdots<i_n\}\), write
\(\epsilon_I=\epsilon_{i_1}\wedge\cdots\wedge\epsilon_{i_n}\).
Define a complex of free left \(W\)-modules by
\begin{equation}\label{mbh:eq:free-koszul}
 \begin{aligned}
 K^W_n&=W\otimes_S\Lambda^n E_S,\\
 d_K(w\epsilon_I)
   &=\sum_{a=1}^n(-1)^{a-1}
       wp_{i_a}\epsilon_{I\setminus\{i_a\}},&
 \pi_K(w)&=w1\in V_S.
 \end{aligned}
\end{equation}
Commutation of the \(p_m\) gives \(d_K^2=0\).
For comparison, write
\[
 P^W_n=W\otimes_S W^{\otimes_S n}\otimes_S V_S
\]
for the unnormalized left bar resolution.  Its differential \(d_P\)
has the faces of Theorem~\ref{mb:thm:mode-bar}, with the first
coefficient in the regular left module \(W\), and its augmentation is
\(\pi_P(a[]v)=av\).  It is free as a left \(W\)-module on the basis
\begin{equation}\label{mbh:eq:bar-basis}
 e=1[a_1|\cdots|a_n]q^\alpha,
\end{equation}
where every \(a_j\) is an ordered monomial.

\begin{lemma}[A contraction over the coefficient ring]
\label{mbh:lem:integral-contraction}
For an ordered basis vector \(q^\alpha p^\beta\epsilon_I\),
put \(J=\operatorname{supp}\beta\cup I\).
If \(J=\varnothing\), set \(h_K(q^\alpha)=0\).
Otherwise let \(m=\min J\) and set
\begin{equation}\label{mbh:eq:least-index}
 h_K(q^\alpha p^\beta\epsilon_I)=
 \begin{cases}
 0,&m\in I,\\
 q^\alpha p^{\beta-e_m}\epsilon_m\wedge\epsilon_I,&m\notin I.
 \end{cases}
\end{equation}
Here \(e_m\) is the multi-index with \(1\) in position \(m\)
and \(0\) elsewhere.  Extend this map by \(S\)-linearity, and let
\(i_K:V_S\to K^W_0\) be the polynomial inclusion.  Then
\begin{equation}\label{mbh:eq:koszul-contraction}
 d_Kh_K+h_Kd_K=1-i_K\pi_K,\qquad h_K^2=0,
\end{equation}
where \(i_K\pi_K\) acts only in degree zero.  In particular,
\(K^W_\bullet\to V_S\) is a free resolution.
\end{lemma}

\begin{proof}
If \(J=\varnothing\), the input is a polynomial in the \(q_m\)
in degree zero, and both sides of the first identity vanish.
If \(m\in I\), the homotopy vanishes on the input.  In its
differential, only the term removing \(\epsilon_m\) has nonzero
image under \(h_K\).  Since \(m\) is the first exterior index,
this term has positive sign and recovers the input.

If \(m\notin I\), then \(\beta_m>0\) and \(m<i\) for every
\(i\in I\).  Thus the exterior word in
\eqref{mbh:eq:least-index} is already ordered.  The differential
term removing its initial \(\epsilon_m\) recovers the input.
Every other term cancels the term obtained by first removing the
same exterior generator and then applying \(h_K\): the insertion
of the initial generator changes its sign once.  This proves the
first identity.  After a nonzero application of \(h_K\), the least
occupied index belongs to the exterior set, so a second application
is zero.  Taking \(h_{K,-1}=i_K\) contracts the augmented complex.
\end{proof}

The bar contraction from the proof of
Theorem~\ref{mb:thm:mode-bar} also works over \(S\):
\begin{equation}\label{mbh:eq:bar-contraction}
 h_{P,-1}(v)=1[]v,\qquad
 h_{P,n}\bigl(a_0[a_1|\cdots|a_n]v\bigr)
       =1[a_0|a_1|\cdots|a_n]v.
\end{equation}
Its first face is the input and the other faces cancel those in
\(h_Pd_P\), so \(d_Ph_P+h_Pd_P=1\) on the augmented complex.
Both \(h_K\) and \(h_P\) are only \(S\)-linear.
For example, \(h_K(p_m)=\epsilon_m\), whereas
\(p_mh_K(1)=0\).  The module-linear comparison must therefore be
defined on free bases and then extended by \(W\)-linearity.

Define \(F:K^W\to P^W\) by antisymmetrization:
\begin{equation}\label{mbh:eq:antisymmetrization}
 F_n(w\epsilon_I)=
 \sum_{\sigma\in\mathfrak S_n}\operatorname{sgn}(\sigma)\,
 w[p_{i_{\sigma(1)}}|\cdots|p_{i_{\sigma(n)}}]1,
 \qquad F_0(w)=w[]1.
\end{equation}
Here \(\mathfrak S_n\) is the permutation group of \(\{1,\ldots,n\}\).
For the reverse map, prescribe
\(G_0(1[]q^\alpha)=q^\alpha\), and, for \(n>0\), prescribe
on each basis element \(e\) of \eqref{mbh:eq:bar-basis}
\begin{equation}\label{mbh:eq:reverse-recursion}
 G_n(e)=h_{K,n-1}\bigl(G_{n-1}(d_Pe)\bigr).
\end{equation}
Expand the argument in ordered monomials before applying \(h_K\),
and extend each \(G_n\) by left \(W\)-linearity.
Finally put \(D=1_{P^W}-FG\), set \(H_{-1}=0\), and define
on the same free basis
\begin{equation}\label{mbh:eq:bar-homotopy-recursion}
 H_n(e)=h_{P,n}\bigl(D_n(e)-H_{n-1}d_Pe\bigr),
 \qquad n\ge0,
\end{equation}
with the second term omitted for \(n=0\).
Again extend by left \(W\)-linearity.  These prescriptions apply
to the displayed free basis; they do not apply the \(S\)-linear
contractions to an arbitrary left multiple of a basis element.

\begin{theorem}[An integral retraction of the oscillator bar]
\label{mbh:thm:retraction}
The maps \(F\) and \(G\) are left \(W\)-linear chain maps lifting
the identity of \(V_S\), and
\begin{equation}\label{mbh:eq:retraction-identities}
 GF=1_{K^W},\qquad d_PH+Hd_P=1_{P^W}-FG.
\end{equation}
All formulas terminate on each element, introduce no new oscillator
index, and commute with every base change \(S\to S'\), with each
\(c_m\) replaced by its image in \(S'\).  They use no division.
\end{theorem}

\begin{proof}
For \(F\), the first bar faces give \(Fd_K\).
The last faces vanish since \(p_m1=0\), and the interior faces
cancel in pairs by interchanging two adjacent commuting \(p_m\).
The signs in each pair are opposite, also in characteristic two.
Thus \(d_PF=Fd_K\), and the augmentations agree.

The identity \(\pi_KG_0=\pi_P\) follows on the free basis.
Assume the chain identity for \(G\) in lower degrees, and let
\(e\in P^W_n\) be a free basis element.  Put
\(z=G_{n-1}d_Pe\).  If \(n=1\), then
\(\pi_Kz=\pi_Pd_Pe=0\).  If \(n>1\), then
\(d_Kz=G_{n-2}d_P^2e=0\).  In either case,
\eqref{mbh:eq:koszul-contraction} gives \(d_Kh_Kz=z\).
This proves \(d_KG_n(e)=G_{n-1}d_Pe\).
Both sides are \(W\)-linear, so the identity holds everywhere.

In degree zero, \(G_0F_0=1\).  Suppose \(GF=1\) in degree
\(n-1\), and take \(n>0\).
The vector \(F_n(\epsilon_I)\) is an \(S\)-linear combination
of free basis elements of \(P^W_n\).  Applying
\eqref{mbh:eq:reverse-recursion} term by term therefore gives
\[
 G_nF_n(\epsilon_I)
   =h_KG_{n-1}d_PF_n(\epsilon_I)
   =h_Kd_K(\epsilon_I)=\epsilon_I.
\]
For the last equality, \(h_K(\epsilon_I)=0\) and
\eqref{mbh:eq:koszul-contraction} applies in positive degree.
Left \(W\)-linearity proves the first identity in
\eqref{mbh:eq:retraction-identities}.

For the second identity, \(\pi_PD_0=0\), since \(F\) and \(G\)
lift the identity of \(V_S\).  Thus
\(d_PH_0(e)=D_0(e)\) by \eqref{mbh:eq:bar-contraction}.
Assume the homotopy identity below degree \(n\), and put
\(z=D_n(e)-H_{n-1}d_Pe\).
Since \(D\) is a chain map, the preceding degree gives
\[
 d_Pz
  =D_{n-1}d_Pe
    -(D_{n-1}-H_{n-2}d_P)d_Pe=0.
\]
The bar contraction yields \(d_Ph_Pz=z\), proving the required
identity on the free basis and hence on all left multiples.

Each recursion lowers homological degree.  Reordering a finite
product gives finitely many ordered monomials, and none of the
contractions introduces an index absent from the input.  Thus the
maps are compatible with inclusions of finite sets of modes, with
their order inherited from the positive integers.
Their coefficients use only addition, multiplication and signs in
the \(c_m\).  The same free-basis formulas after any scalar change
satisfy the same identities; no flatness assumption on that change
is required.
\end{proof}

In particular, antisymmetrization does not require division by a
factorial.  For \(i<j\), the recursion gives
\[
 \begin{aligned}
 G_1(1[p_i]1)&=\epsilon_i,&
 G_1(1[p_ip_j]1)&=p_j\epsilon_i,\\
 G_2(1[p_i|p_j]1)&=\epsilon_i\wedge\epsilon_j,&
 G_2(1[p_j|p_i]1)&=0.
 \end{aligned}
\]
Their antisymmetric difference maps to
\(\epsilon_i\wedge\epsilon_j\), including in characteristic two.

For any right \(W\)-module \(L\), tensoring the maps and homotopy
over \(W\) gives a chain homotopy equivalence
\begin{equation}\label{mbh:eq:module-comparison}
 B_S(L,W,V_S)
 \ \underset{F_L}{\stackrel{G_L}{\rightleftarrows}}\
 K(L),\qquad
 K(L)_n=L\otimes_S\Lambda^n E_S,
\end{equation}
where \(F_L=1_L\otimes_WF\) goes from \(K(L)\) to the bar,
\(G_L=1_L\otimes_WG\) goes from the bar to \(K(L)\), and
\begin{equation}\label{mbh:eq:annihilator-differential}
 d(\lambda\epsilon_I)=
 \sum_{a=1}^n(-1)^{a-1}
       (\lambda p_{i_a})\epsilon_{I\setminus\{i_a\}}.
\end{equation}
The identities \(G_LF_L=1\) and
\(dH_L+H_Ld=1-F_LG_L\), with \(H_L=1_L\otimes_WH\),
follow by tensoring \eqref{mbh:eq:retraction-identities}.
This requires no flatness of \(L\) over \(W\).

We now use the actual restricted-dual end.  Give \(q_m\) weight
\(m\), write \((V_S)_a\) for the weight-\(a\) component, and put
\[
 M_S=\bigoplus_{a\ge0}\operatorname{Hom}_S((V_S)_a,S).
\]
Each weight component is finite free.  Its monomial duals form a
basis \(\lambda_\alpha\), characterized by
\(\lambda_\alpha(q^\beta)=\delta_{\alpha,\beta}\).
The right action \((\lambda a)(v)=\lambda(av)\) gives
\begin{equation}\label{mbh:eq:dual-actions}
 \lambda_\alpha p_m
     =c_m(\alpha_m+1)\lambda_{\alpha+e_m},
 \qquad
 \lambda_\alpha q_m=
 \begin{cases}
 \lambda_{\alpha-e_m},&\alpha_m>0,\\
 0,&\alpha_m=0.
 \end{cases}
\end{equation}
Thus \eqref{mbh:eq:annihilator-differential} computes the
derived tensor product on the stated Fock-module ends over \(S\).
The evaluation of the bar, \(\lambda[]v\mapsto\lambda(v)\)
in degree zero, becomes
\begin{equation}\label{mbh:eq:evaluation}
 \pi(\lambda_\alpha)=\delta_{\alpha,0}
\end{equation}
on \(K(M_S)_0\), with zero value in positive degrees.
For \(F_0\) this follows from \(F_0(w)=w[]1\).
For \(G_0\) it follows from
\((\lambda aq^\alpha)(1)=\lambda(aq^\alpha1)\).
The map is a chain map because \((\lambda p_m)(1)=0\).

\begin{theorem}[Evaluation at nonzero level]
\label{mbh:thm:nonzero}
Let \(k\) be a field of characteristic zero and \(c\in k^\times\).
For \(\mathscr W_c\), \(V=k[q_1,q_2,\ldots]\) and its
restricted dual \(M\), the evaluation
\[
 B_k(M,\mathscr W_c,V)\longrightarrow k
\]
is a quasi-isomorphism.  Equivalently,
\[
 \operatorname{Tor}^{\mathscr W_c}_j(M,V)=
 \begin{cases}k,&j=0,\\0,&j>0.\end{cases}
\]
\end{theorem}

\begin{proof}
All factors in
\begin{equation}\label{mbh:eq:fixed-normalization}
 k[x_1,x_2,\ldots]\longrightarrow M,\qquad
 x^\alpha\longmapsto
 \left(\prod_m(cm)^{\alpha_m}\alpha_m!\right)\lambda_\alpha
\end{equation}
are invertible.  This is a vector-space isomorphism, and
\eqref{mbh:eq:dual-actions} identifies right multiplication by
\(p_m\) with multiplication by \(x_m\).
The complex \(K(M)\) becomes
\[
 K=k[x_1,x_2,\ldots]\otimes_k\Lambda(\epsilon_m:m\ge1),
 \qquad
 \delta=\sum_m x_m\iota_m,
\]
where \(\iota_m\) is exterior contraction:
\(\iota_m(\epsilon_I)=(-1)^{a-1}
\epsilon_{I\setminus\{m\}}\) when \(m=i_a\), and zero
when \(m\notin I\).

Give both \(x_m\) and \(\epsilon_m\) auxiliary weight one.
This is separate from the state weight and homological degree.
On auxiliary weight \(w>0\), define
\begin{equation}\label{mbh:eq:euler-homotopy}
 h=\frac1w\sum_m\epsilon_m\wedge\partial_{x_m},
 \qquad h(1)=0.
\end{equation}
The relations
\(\iota_i(\epsilon_j\wedge-)+\epsilon_j\wedge\iota_i
=\delta_{ij}\) and the polynomial product rule give
\[
 \delta\Bigl(\sum_m\epsilon_m\wedge\partial_{x_m}\Bigr)
 +\Bigl(\sum_m\epsilon_m\wedge\partial_{x_m}\Bigr)\delta
 =\sum_m x_m\partial_{x_m}+\sum_m\epsilon_m\wedge\iota_m.
\]
The last operator is multiplication by \(w\).
Thus \(\delta h+h\delta=1-i\pi\), where \(i(1)=1\) and
\(\pi\) is the constant coefficient in degree zero.
Every sum is finite on its input.  This contracts \(K\) onto
\(k\), with exactly the evaluation \eqref{mbh:eq:evaluation}.
The equivalence \eqref{mbh:eq:module-comparison} proves the claim.
\end{proof}

This scalar homology uses the two module actions.
It supplies no character of \(\mathscr W_c\), which remains
excluded by \([p_1,q_1]=c1\).
Over a parameter ring, division by the level would lose homology.

\begin{theorem}[The level family and derived specialization]
\label{mbh:thm:family}
Let \(k\) be a field of characteristic zero, put \(S=k[\kappa]\),
and take \(c_m=\kappa m\).  The bar on \(V_S\) and \(M_S\)
is chain homotopy equivalent to
\begin{equation}\label{mbh:eq:family-complex}
 K_S=S[x_1,x_2,\ldots]\otimes_S\Lambda(\epsilon_m:m\ge1),
 \qquad d_\kappa=\kappa\delta,\qquad
 \delta=\sum_m x_m\iota_m.
\end{equation}
Set \(B_j=\operatorname{im}(\delta:K_{S,j+1}\to K_{S,j})\)
and \(B_{-1}=0\).  Then
\begin{equation}\label{mbh:eq:family-homology}
 H_j(K_S)=
 \begin{cases}
 S\oplus B_0/\kappa B_0,&j=0,\\
 B_j/\kappa B_j,&j>0,
 \end{cases}
 \qquad
 B_0=(x_1,x_2,\ldots)S[x_1,x_2,\ldots].
\end{equation}
The kernel of evaluation on homology,
\(H_\bullet(K_S)\to S\) with \(S\) in degree zero, is
annihilated by \(\kappa\).  It is nonzero in every positive
homological degree.
The derived zero fibre is
\begin{equation}\label{mbh:eq:derived-zero}
 H_j(K_S\otimes_S^{\mathbb L}k)
   =k[x_1,x_2,\ldots]\otimes_k\Lambda^j(\epsilon_m:m\ge1),
 \qquad k=S/(\kappa).
\end{equation}
\end{theorem}

\begin{proof}
Normalize the monomial duals by
\[
 x^\alpha\longmapsto
 \left(\prod_m m^{\alpha_m}\alpha_m!\right)\lambda_\alpha.
\]
The factors are units in \(k\), so this is an invertible
\(S\)-linear change of basis.  It introduces no inverse of
\(\kappa\).  Formula~\eqref{mbh:eq:dual-actions} then gives
\eqref{mbh:eq:family-complex}.
The homotopy \eqref{mbh:eq:euler-homotopy} is \(S\)-linear and
satisfies
\begin{equation}\label{mbh:eq:family-homotopy}
 d_\kappa h+hd_\kappa=\kappa(1-i\pi).
\end{equation}
Here \(\pi\) takes the constant coefficient in the variables
\(x_m\), with values in \(S\), and \(i\) is the scalar inclusion.

Let \(K_{S,j}^+\) denote the part of positive auxiliary weight;
only the scalar summand \(S\cdot1\) in degree zero is removed.
The identity \(\delta h+h\delta=1\) on this part gives
\begin{equation}\label{mbh:eq:family-splitting}
 K_{S,j}^+=B_j\oplus hB_{j-1},\qquad
 d_\kappa(hb)=\kappa b,\qquad d_\kappa(B_j)=0.
\end{equation}
Indeed, every \(z\) is
\(\delta hz+h\delta z\), which lies in the stated sum.
If \(hb\in B_j\) with \(b\in B_{j-1}\), applying
\(\delta\) gives \(b=0\), proving directness.
Since \(\delta\) and \(h\) are defined over \(k\), each \(B_j\)
is \(S\) tensored with a \(k\)-vector space and hence is
\(S\)-free.  Multiplication by \(\kappa\) is injective on it.
The kernels and images in \eqref{mbh:eq:family-splitting}
now give \eqref{mbh:eq:family-homology}.

For \(j>0\), the cycle
\begin{equation}\label{mbh:eq:torsion-cycle}
 z_j=\delta(\epsilon_1\wedge\cdots\wedge\epsilon_{j+1})
   =\sum_{a=1}^{j+1}(-1)^{a-1}x_a\,
      \epsilon_1\wedge\cdots\wedge
      \widehat{\epsilon_a}\wedge\cdots\wedge\epsilon_{j+1}
\end{equation}
lies in \(B_j\) and has nonzero coefficients independent of
\(\kappa\).  It cannot lie in \(\kappa B_j\), so its class is
nonzero.  The equality
\(d_\kappa(\epsilon_1\wedge\cdots\wedge\epsilon_{j+1})
=\kappa z_j\) also displays its annihilation.

The splitting \eqref{mbh:eq:family-splitting} expresses \(K_S\)
as the scalar \(S\) and a direct sum of two-term free complexes
\[
 hB_j\xrightarrow{\ \kappa\ }B_j,
\]
in degrees \(j+1\) and \(j\).  A two-term complex of free
modules is \(K\)-flat: tensoring it with an acyclic complex is
the cone of a map between acyclic direct sums of that complex.
Direct sums preserve exactness.  The same is therefore true of
\(K_S\), so its ordinary tensor product with \(k=S/(\kappa)\)
computes the derived tensor product.  Its differential becomes zero,
proving \eqref{mbh:eq:derived-zero}.
\end{proof}

The extra specialization classes can be identified without a
convergence argument.  For \(j\ge1\), the preceding splitting gives
the exact sequence
\begin{equation}\label{mbh:eq:base-change}
 \begin{aligned}
 0\longrightarrow H_j(K_S)\otimes_S k
 &\longrightarrow H_j(K_S\otimes_S^{\mathbb L}k)\\
 &\longrightarrow
 \operatorname{Tor}_1^S(H_{j-1}(K_S),k)
 \longrightarrow0.
 \end{aligned}
\end{equation}
Its first term is represented by \(B_j\otimes_Sk\), and its
last term by \(hB_{j-1}\otimes_Sk\), using
\(h:B_{j-1}\to hB_{j-1}\).
Indeed, \(0\to B_{j-1}\xrightarrow{\kappa}B_{j-1}
\to B_{j-1}/\kappa B_{j-1}\to0\) is a free resolution,
and the scalar summand of \(H_0\) has zero \(\operatorname{Tor}_1\).
For a single oscillator, \(B_1=0\): the family has \(H_1=0\),
whereas its zero fibre has
\(H_1=k[x_1]\epsilon_1\).  These classes come entirely from
the last term of \eqref{mbh:eq:base-change}.
With infinitely many oscillators,
\eqref{mbh:eq:torsion-cycle} gives family torsion in every
positive degree.  Inverting \(\kappa\) instead gives the
contraction \(\kappa^{-1}h\) and recovers
Theorem~\ref{mbh:thm:nonzero}.

At zero level over any field, every right \(p_m\)-action on
the restricted dual vanishes.  Thus
\eqref{mbh:eq:annihilator-differential} gives
\begin{equation}\label{mbh:eq:zero-level}
 \operatorname{Tor}^{\mathscr W_0}_j(M,V)
   =M\otimes_k\Lambda^j(\epsilon_m:m\ge1).
\end{equation}
The algebra \(\mathscr W_0\) still contains every \(p_m\);
its zero action on these module ends does not remove the
corresponding exterior generator.

\begin{theorem}[The positive-characteristic basis]
\label{mbh:thm:positive-characteristic}
Let \(k\) be a field of characteristic \(\ell>0\), let \(c\in k^\times\),
and use the polynomial Fock module and its restricted dual.
Put
\[
 G=\{m\ge1:\ell\nmid m\},\qquad
 Z=\{m\ge1:\ell\mid m\}.
\]
For every \(j\ge0\), a basis of
\(\operatorname{Tor}^{\mathscr W_c}_j(M,V)\) is represented
in \(K(M)\) by
\begin{equation}\label{mbh:eq:characteristic-basis}
 \lambda_{\alpha(b,I)}\epsilon_I,\qquad
 |I|=j,\qquad
 \alpha(b,I)_m=
 \begin{cases}
 \ell b_m+(\ell-1)\mathbf1_{m\in I},&m\in G,\\
 b_m,&m\in Z,
 \end{cases}
\end{equation}
where \(I\) is a finite set and \(b=(b_m)\) is a finite-support
multi-index of nonnegative integers.
This is a basis of graded vector spaces; no multiplication on
these Tor groups is asserted.
\end{theorem}

\begin{proof}
For a single mode \(m\), let
\(M_m=\bigoplus_{n\ge0}k\lambda_n\).
Its factor in the annihilator complex is
\begin{equation}\label{mbh:eq:one-mode}
 M_m\epsilon_m\longrightarrow M_m,\qquad
 \lambda_n\epsilon_m\longmapsto cm(n+1)\lambda_{n+1}.
\end{equation}
If \(m\in G\), its degree-zero cokernel has basis
\(\lambda_{\ell b}\), \(b\ge0\), and its degree-one kernel
has basis \(\lambda_{\ell b+\ell-1}\epsilon_m\).
Every other degree-zero basis vector is the image of exactly one
basis vector with nonzero coefficient, and every remaining
degree-one vector belongs to such a pair.
If \(m\in Z\), the differential is zero and both copies of
\(M_m\) survive.

The full complex is the algebraic restricted tensor product of
\eqref{mbh:eq:one-mode} over all \(m\), using \(\lambda_0\)
in degree zero at each new index.  For a finite set of indices,
split each factor into the displayed kernel and cokernel
representatives and its two-term pairs with nonzero differential.
Those pairs are contractible over the field \(k\).  Tensoring their
contractions with the other factors proves that the homology of the
finite tensor product is the tensor product of the displayed
homologies.  This splitting preserves the distinguished
degree-zero vector \(\lambda_0\), so it is compatible with
inclusions of finite index sets.  Their union is the full
algebraic complex, and filtered unions of vector spaces preserve
exactness.  Selecting the degree-one factor exactly at the indices
of \(I\) gives precisely \eqref{mbh:eq:characteristic-basis}.
\end{proof}

For example,
\[
 [\,\lambda_{\ell-1}[p_1]1\,]
   \in H_1\bigl(B_k(M,\mathscr W_c,V)\bigr)
\]
is nonzero.  Both bar faces vanish: the first has coefficient
\(c\ell=0\), and the last contains \(p_1 1=0\).
Under \(F_M\), it is the image of
\(\lambda_{\ell-1}\epsilon_1\), whose class is nonzero by
Theorem~\ref{mbh:thm:positive-characteristic};
\(G_MF_M=1\) proves that it is not a boundary.
Thus the evaluation to \(k\) fails to be a quasi-isomorphism
in every positive characteristic, including at nonzero level.
Modes in \(Z\) supply further independent classes.
Changing the polynomial Fock module to a divided-power module
would change the input action in \eqref{mbh:eq:dual-actions}.

These evaluations retain the two module ends and concern ordinary
algebraic complexes.  The balanced cuts of
\eqref{mb:eq:module-cuts} retain their stated target; the
retraction \eqref{mbh:eq:module-comparison} makes no assertion
of a coproduct equivalence.  A geometric collision comparison
or an extension to completed states requires additional maps.


\subsubsection{Coefficient comparison with the same module ends}

The coefficient comparison also has a two-sided version.
Work in characteristic zero and extend scalars to \(R\).
Put \(W=R\otimes_k\mathscr W_\kappa\),
\(M_R=R\otimes_kM\), and \(V_R=R\otimes_kV\).
The algebras \(C\otimes_R W\) and \(A\otimes_R W\)
act on their module ends through their first-endpoint augmentations
and the stated mode actions.  No augmentation of \(W\) is used.

\begin{theorem}[The state-valued two-sided coefficient comparison]
\label{mb:thm:two-sided-comparison}
Assume \(\operatorname{char}k=0\).  The consecutive-partition formula, with
\[
 [g_1ds\otimes a_1|\cdots|g_nds\otimes a_n]
 \longmapsto
 [J_n(g_1,\ldots,g_n)\theta\otimes a_1\cdots a_n]
\]
for its higher components, defines a chain quasi-isomorphism
\[
 B_R(M_R,C\otimes_R W,V_R)
    \longrightarrow B_R(M_R,A\otimes_R W,V_R).
\]
Its linear component preserves every mode and its action on either
Fock-module end.  It preserves the total spatial pole bound, and its
higher components use only consecutive middle blocks.
\end{theorem}

\begin{proof}
For a differential graded algebra \(E\) acting on these degree-zero
module ends, its two-sided bar has middle suspension signs
\eqref{sb:eq:bar-signs}.  Its left outer face has positive sign.
Its right outer face has sign
\((-1)^{\sum_j(|a_j|-1)}\).  Nonzero end actions require
degree-zero coefficients, so these conventions reduce exactly to
Theorem~\ref{mb:thm:mode-bar} when all middle entries have degree zero.
The algebra and module product rules verify its square in the same
way as Lemma~\ref{sb:lem:bar} and Theorem~\ref{mb:thm:mode-bar}.

The internal chain calculation in Lemma~\ref{sb:lem:state-extension}
uses only the differential and multiplication of the state algebra;
it does not use its augmentation.  It therefore applies to \(W\).
For the outer faces, a higher block has coefficient \(\theta\)
and acts by zero at the first endpoint.  Only a singleton block can
act on a module end.  Its action is unchanged by \(I\), because
\(\varepsilon_A I=\varepsilon_C\).
Thus the outer terms agree as well, proving the chain equation.

The filtration by middle word length has associated map
\(\id_{M_R}\otimes(s(I\otimes\id_W))^{\otimes_R r}
 \otimes\id_{V_R}\).
The explicit coefficient contraction proves that this is a chain
homotopy equivalence, even after these tensor products.
Finite filtration induction and passage to the union prove the
quasi-isomorphism.  Products of modes do not introduce spatial poles,
so the proof of the pole bound is unchanged.
\end{proof}

This bar retains all the finite mode words that define the fields.
It does not authorize substitution of a collision variable into an
arbitrary formal distribution of mode words.  For example, the formal
middle entry \(\sum_{m\ge1}(q_mz^{m-1}+p_mz^{-m-1})\) has two unbounded
tails in the mode algebra, even when sufficiently large annihilation
modes kill a fixed right-end state.  Its matrix coefficients and the
specific rational collision representatives must be expanded in their
stated domains before a diagonal substitution is made.


\subsection{Induction over a fixed binary tree}

In characteristic zero the coefficient comparison extends through any prescribed finite
sequence of binary order intervals.  We state its domain and the
extent of its coherence explicitly.

Let a binary tree have internal vertices \(v_1,\ldots,v_m\),
in a chosen processing order.  At vertex \(v\), choose a local
child difference \(u_v\) and form the two algebras \(C_v,A_v\)
of \eqref{sb:eq:two-chart} with \(u,s,\theta\) replaced by
\(u_v,s_v,\theta_v\).
For an augmented differential graded state algebra \(D\) over \(k\),
put \(R_T=k[u_v:v\text{ internal}]\) and define
\[
 D_T^C=\left(\bigotimes_{v}C_v\right)\otimes D,
 \qquad
 D_T^A=\left(\bigotimes_{v}A_v\right)\otimes D.
\]
Tensor products are graded tensor products of algebras; their order is fixed.
Their augmentations take the first endpoint at every vertex and the
state augmentation in \(D\), with values in \(R_T\).
All bars in this section are relative to \(R_T\).  Each local
coefficient algebra is first extended from \(k[u_v]\) to \(R_T\).
At step \(j\), move the factor \(C_{v_j}\) to the first
position by the graded tensor permutation, apply
Lemma~\ref{sb:lem:state-extension} with all other factors as its state
algebra, and move \(A_{v_j}\) back.  Compose these coalgebra maps
for \(j=1,\ldots,m\).  Call the result \(\mathcal I_T\).

\begin{theorem}[Finite-tree coefficient bars]\label{sb:thm:tree-bars}
In characteristic zero this recursion gives an explicit coalgebra
chain quasi-isomorphism
\[
 \mathcal I_T:B_{R_T}(D_T^C)\longrightarrow B_{R_T}(D_T^A).
\]
It preserves cuts and the total spatial pole tuple in every word length.
Every component is a finite sum of iterated polynomial integrals and
products.  It commutes with spatial substitutions and formal collision
expansions on the fixed-input domains described below.
The assertion is for the specified processing order; it does not assert
independence of that order or compatibility with every tree rotation.
\end{theorem}

\begin{proof}
Each step is a coalgebra chain quasi-isomorphism by
Lemma~\ref{sb:lem:state-extension}; the graded tensor permutations are
strict algebra isomorphisms.  Their finite composite has the same
properties.  Each higher component groups consecutive input words,
so the recursion never increases word length and has finitely many
terms on a fixed word.  A spatial coefficient of an original entry
appears exactly once in each final product, so its pole contribution
is added once.  This proves the bound by induction on the processed
vertices.  The same recursion proves the cut identity, since it holds
at each step.

For a collision expansion, place one child of each vertex at coordinate
zero and the other at \(u_v\).  Attach a scale \(\epsilon_v\)
to each nonroot internal vertex.  A leaf coordinate is the sum of
local displacements multiplied by the scales above them.  At the
lowest common ancestor of two leaves their difference is
\[
 z_i-z_j=\lambda_v(u_v+h),
\]
up to a fixed sign, where \(\lambda_v\) is an ancestor-scale
monomial and \(h\) has positive descendant-scale degree.
On an overlap where the comparison varies, \(u_v\) is invertible.
Its rational inverse expands geometrically after the finite factor
\(\lambda_v^{-1}\).

For fixed polynomial inputs and finitely many original denominators,
these expansions lie in
\[
 B_T[[\epsilon_v]][(\prod_v\epsilon_v)^{-1}],
\]
where \(B_T\) is the specified tensor product of coordinate overlap
rings.  First-endpoint evaluation commutes with these substitutions,
since they change spatial coordinates only.  Thus the relative
augmentations agree after extension to the formal spatial base.
The negative exponent of each scale has a common finite lower
bound.  A fixed scale coefficient and state weight uses finitely many
terms.  Order-coordinate integration and each finite consecutive
partition therefore commute with coefficient extraction.  Neither
operation changes spatial variables.  This proves the asserted
naturality on these domains.
\end{proof}

When a prescribed grafting procedure orders the new vertices, the
recursion applies to that procedure and supplies all bar components
at every step.  It supplies actual higher corrections to coefficient
multiplication.  An operadic comparison would additionally identify
different processing orders and tree rotations by homotopies that
respect their boundary faces.  Such data are not implied by composing
chain quasi-isomorphisms.

The three-leaf comparison in Theorem~\ref{tr:thm:bar-rotation} constructs
such homotopies on a specified localization, with its polynomial-form
comparison in Proposition~\ref{tr:prop:form-square}.  That localization
retains one collision divisor.  Proposition~\ref{tr:prop:global-obstruction}
shows why the same equivalence cannot hold for the two whole
coefficient complexes over their original spatial ring.  A comparison
retaining every boundary and compatible with successive rotations
requires further constructions.

The completion restriction is necessary.  The series
\[
 \sum_{n\ge0}\epsilon^n u^{-n^2}
\]
is valid in \(k[u,u^{-1}][[\epsilon]]\), but replacing \(u\)
by a second collision scale gives no lower bound in that scale.
It does not lie in the simultaneous Laurent ring just displayed.
The state operations have another restriction: in the free state
completion, \(\eta_r\to0\), while the scalar component of
\(Y_A(x_0,\eta_r)\) is \(\theta/u^{r+1}\), which does not
tend to zero in its discrete weight-zero coefficient space.
Neither arbitrary intermediate substitutions nor arbitrary completed
state inputs are covered by the finite-tree construction.

\subsection{From coefficient bars to geometric collisions}
\label{sb:sec:state-native-boundary}

The proved comparison has a precise boundary.  Polynomial-form
coefficients and ordered overlap coefficients have quasi-isomorphic
bars in characteristic zero, with every higher component explicitly
given.  Their binary free-state operations agree under the linear
component.  In positive characteristic the direct overlap operation
retains the pole class, while any comparison from the same polynomial-form
complex preserving it is impossible.

The native bar with two unit ends requires further data:
an actual augmented algebra object whose multiplication and augmentation
are collision-preserving maps.  The normal product and vacuum projection
of nonzero-level Heisenberg states fail these requirements by the
computed length-three associator and binary pole class.
A comparison retaining those Heisenberg collisions must therefore
change the proposed augmented carrier or change the type of bar.
Such a change must include an explicit map on states and coefficients;
it cannot be achieved by reinterpreting the displayed constant product.
The mode bar supplies one such change of algebraic carrier: it uses
the actual Fock-module actions and computes their derived tensor product,
with the stated balanced cuts.  The integral retraction of
Theorem~\ref{mbh:thm:retraction} evaluates that tensor product on the
same module ends.  For a characteristic-zero field \(k\), its
evaluation at nonzero level is a quasi-isomorphism to \(k\).
Over \(k[\kappa]\), the kernel of evaluation on homology is
annihilated by \(\kappa\).  Positive characteristic has the
explicit classes of
Theorem~\ref{mbh:thm:positive-characteristic}.
The two-sided coefficient comparison preserves every finite mode word
in characteristic zero.  Neither comparison constructs a map to the
geometric collision carrier or identifies the balanced cuts with a
coproduct on a scalar-ended bar.

\begin{problem}[A collision-preserving bar carrier]
Construct a geometric chain carrier for Heisenberg collisions whose
faces retain the singular coefficients and their central vacuum term.
Specify its unit or augmentation data, collision pullbacks, state
realization, and total grading.  Give a chain map from a stated
coefficient presentation, and compare the two arity-three collision
domains before constructing homotopies for their higher refinements.
If an internal augmented associative bar is used, identify the added
states or changed target that evade
Proposition~\ref{sb:prop:augmentation-obstruction}; if another type of
bar is used, define its differential and cuts directly.
\end{problem}
